% =====================================================================
% Platonic synthesis after the twenty-agent adversarial wave.
% Author: Raeez Lorgat.  Chriss--Ginzburg register; subscripted
% $\kappa_\bullet$ throughout; no AI attribution; no bookkeeping
% vocabulary in reader-facing prose.
%
% Inscription target: the central spine of working_notes.tex, immediately
% after the abstract.  Each theorem here is a survivor of the adversarial
% pass; each retraction states the wrong claim, the precise error, and
% the ghost-theorem that survives.  Proofs reach CFG-2026 detail where
% the chain level is the correct lane; the $(\infty,1)$-categorical
% lane is labelled where it is primary.
% =====================================================================

\providecommand{\ClaimStatusTheorem}{\textsc{[Theorem]}}
\providecommand{\ClaimStatusConjectured}{\textsc{[Conjectured]}}
\providecommand{\ClaimStatusCorrected}{\textsc{[Corrected]}}
\providecommand{\ClaimStatusOpen}{\textsc{[Open]}}
\providecommand{\ClaimStatusRetracted}{\textsc{[Retracted]}}
\providecommand{\ClaimStatusDefinition}{\textsc{[Definition]}}
\providecommand{\kBKM}{\kappa_{\mathrm{BKM}}}
\providecommand{\kch}{\kappa_{\mathrm{ch}}}
\providecommand{\kcat}{\kappa_{\mathrm{cat}}}
\providecommand{\kfib}{\kappa_{\mathrm{fiber}}}
\providecommand{\kanom}{\kappa_{\mathrm{anom}}}
\providecommand{\ZZ}{\mathbb{Z}}
\providecommand{\CC}{\mathbb{C}}
\providecommand{\QQ}{\mathbb{Q}}
\providecommand{\CYcat}{\mathrm{CY}\text{-cat}}
\providecommand{\ChirAlg}{\mathrm{ChirAlg}}
\providecommand{\HolFA}{\mathrm{HolFA}}
\providecommand{\hCS}{\mathrm{hCS}}
\providecommand{\Obs}{\mathrm{Obs}}
\providecommand{\CEch}{\mathrm{CE}^{\bullet}_{\bar\partial,\mathrm{chir}}}
\providecommand{\Ufact}{\mathbf{U}^{\mathrm{fact}}}
\providecommand{\cF}{\mathcal{F}}
\providecommand{\cE}{\mathcal{E}}
\providecommand{\cA}{\mathcal{A}}
\providecommand{\cO}{\mathcal{O}}
\providecommand{\cL}{\mathcal{L}}
\providecommand{\cH}{\mathcal{H}}
\providecommand{\cW}{\mathcal{W}}
\providecommand{\cM}{\mathcal{M}}
\providecommand{\cV}{\mathcal{V}}
\providecommand{\fg}{\mathfrak{g}}
\providecommand{\Yplus}{Y^+_{\epsilon_1, \epsilon_2, \epsilon_3}}
\providecommand{\CoHA}{\mathrm{CoHA}}
\providecommand{\Winf}{\mathcal{W}_{1+\infty}}
\providecommand{\Mtf}{M_{24}}
\providecommand{\bH}{\mathbf{H}}
\providecommand{\FMcpt}{\mathrm{FM}_{\CC^3}}
\providecommand{\GRTone}{\mathrm{GRT}_1}
\providecommand{\Mpr}{\mathrm{Mp}}
\providecommand{\Spr}{\mathrm{Sp}}
\providecommand{\SOr}{\mathrm{SO}_+}
\providecommand{\Mukr}{\mathrm{Muk}}
\providecommand{\Wttl}{\widetilde{\Lambda}}
\providecommand{\LLeech}{\Lambda_{\mathrm{Leech}}}

\section{The surviving platonic synthesis}
\label{wn:sec:platonic-synthesis-spine}

\emph{This section is the spine of the monograph. Every other section
computes a consequence, grounds a witness, or opens a frontier that
this spine organises. The statements below have survived an
adversarial pass that retracted nine prior claims and sharpened the
scope of every theorem; each retracted claim is restated in
\S\ref{wn:subsec:spine-retractions} together with the
ghost-theorem it conceals. Where the chain-level and
$(\infty,1)$-categorical lanes both carry a proof, both are stated
with their native scope; no lane subsumes the other.}

\subsection{The canonical object: one $E_d^{\mathrm{hol}}$-factorisation
algebra, many $E_1$-chiral shadows}
\label{wn:subsec:spine-canonical-object}

\begin{theorem}[Two-stage factorisation of $\Phi_d$, at honest scope]
\label{wn:thm:spine-two-stage}\ClaimStatusTheorem
The Calabi--Yau-to-chiral functor decomposes as
\[
\Phi_d \;=\; \mathrm{Sp}^{\mathrm{ch}}_{\Sigma_{d-1}, C} \;\circ\;
\Phi^{\mathrm{FA}}_d,\qquad
\Phi^{\mathrm{FA}}_d \colon \CYcat_d^{\mathrm{sm,\,prop}} \longrightarrow
E_d^{\mathrm{hol}}\text{-}\HolFA(X),\qquad
\mathrm{Sp}^{\mathrm{ch}}_{\Sigma_{d-1}, C} \colon
E_d^{\mathrm{hol}}\text{-}\HolFA(X) \longrightarrow
E_{n_{\mathrm{nat}}(d)}\text{-}\ChirAlg(C),
\]
with $n_{\mathrm{nat}}(d) = \infty, 2, 1, 1, 1, \dots$ at
$d = 1, 2, 3, 4, 5, \dots$ forced by Dunn--Lurie additivity. The
forcing argument is a one-line calculation the reader should perform:
Dunn--Lurie gives $E_p \otimes E_q \simeq E_{p+q}$ on the topological
side; here one axis is holomorphic (contributes $E_d^{\mathrm{hol}}$,
meaning the full holomorphic multiplication in $d$ complex directions)
and a transverse $(d-1)$-cycle contracts $d-1$ topological directions
to a point (each directional integration removes one $E_1$-factor).
The residue on the reference curve is then
$E_d^{\mathrm{hol}} \otimes E_{-(d-1)}^{\mathrm{top}}$, which is
$E_1$ in the $d\ge 3$ regime where the holomorphic axis survives as
the chiral direction while all transverse topological directions have
been integrated away; at $d=2$ the residue retains one holomorphic
direction and one topological direction uncontracted, giving $E_2$;
at $d=1$ nothing contracts and the shadow is the full $E_\infty$ of
complex-analytic algebras over a point (the $\infty$ entry). This is
the arithmetic behind the $\infty, 2, 1, 1, 1, \dots$ sequence.

The specialisation datum is a \emph{transverse oriented framed
submanifold with tubular neighbourhood}, not a bare homology class.
Why the refinement is needed: a bare homology class is an equivalence
class of cycles modulo boundaries, which does not carry enough
information to specify a factorisation-homology integration. Factorisation
homology of an $E_d^{\mathrm{hol}}$-algebra over a $(d-1)$-manifold
requires an oriented framed transversal embedding so that the product
structure on the algebra has a well-defined way to act on inputs
sitting on the cycle. Two cycles in the same transverse-bordism class
produce canonically isomorphic shadows because a transverse bordism
provides an $E_d^{\mathrm{hol}}$-algebra homotopy between the two
integrations, and the homotopy descends to a canonical isomorphism
of specialised $E_1$-chiral algebras.

\emph{Stage~1 is canonical} in the following precise two-lane sense.
On objects, in the rational chain-level lane for $d \in \{1,2\}$, the
map $\Phi^{\mathrm{FA}}_d$ is an $(\infty,1)$-functor and the space of
$E_d$-enhancements compatible with Costello--Gwilliam--Li locality is
contractible. The mechanism: the little-$d$-discs operad has a rational
Koszul model (the Kontsevich graph operad at $d=2$; Tamarkin's
associator-dependent model at higher $d$); a Koszul model with
contractible enhancement space forces the moduli of $E_d$-structures
compatible with a given $\bar\partial$-differential on the Dolbeault
chain complex to be a contractible simplicial set. Concretely, if one
$E_d$-enhancement exists, every other is connected to it by a path in
the space of coherence operations, and each path is unique up to
further higher-order paths --- this is the content of ``contractible''.
The two named primary sources deliver exactly this: Kontsevich's 1999
proof of $E_2$-formality constructs the graph operad $\mathrm{Graphs}_2$
and proves its quasi-isomorphism with the chain operad of Fulton--MacPherson
configurations; Tamarkin's 2003 extension produces a Drinfeld-associator
parameter family of $E_n$-formality isomorphisms, and the parameter
space is a torsor under the Grothendieck--Teichm\"uller group whose
identity component is contractible.

At $d \geq 3$, on objects the functor exists and is unique up to
quasi-isomorphism; on morphisms its chain-level lift is defined up to
a torsor under
$H^{0}_{\mathrm{Lie}}(\mathfrak{grt}_1;
\mathrm{Aut}(\Phi^{\mathrm{FA}}_d(\cC)))$, where $\mathfrak{grt}_1$ is
the Grothendieck--Teichm\"uller Lie algebra acting on $E_d$-formality
quasi-isomorphisms. Why this specific torsor: the ambiguity in a
chain-level $E_d$-enhancement is precisely the ambiguity in a Drinfeld
associator; associators form a torsor under $\mathfrak{grt}_1$ (Drinfeld
1990); and the morphism lift of $\Phi^{\mathrm{FA}}_d$ records which
associator was used on each side. The torsor trivialises on formal
cyclic $A_\infty$ categories --- because formality collapses the
associator-dependent higher brackets to zero, leaving no ambiguity to
parametrise --- and carries genuine content on non-formal categories
such as local~$\bP^2$, where the Atiyah-class cup-products obstruct
formality and the associator choice affects the $\ell_3, \ell_4, \dots$
brackets observably.

Over $\ZZ$, at $d \geq 3$, the topological $E_d$-operad and the
algebraic $E_d$-operad built from the Gerstenhaber bracket of degree
$1-d$ agree over $\QQ$ but are not known to agree integrally; the
statement ``$\Phi^{\mathrm{FA}}_d$ is canonical up to contractible
choice'' is \emph{open} integrally at $d \geq 3$.
\end{theorem}

\begin{corollary}[Many $E_1$-chiral shadows from one $\cF_X$]
\label{wn:cor:spine-many-shadows}
On a fixed $X$, the collection of $E_1$-chiral algebras produced by
$\Phi_d$ is parametrised by transverse-bordism classes of
$(\Sigma_{d-1}, C)$-pairs. These classes are invariants of the
Stage-$1$ factorisation algebra $\cF_X$, not further applications of
$\Phi_d$. The six routes to the quantum vertex chiral group
$G(K3 \times E)$ are \emph{three} (Stage-$2$) cycle-class
specialisations of $\cF_{K3 \times E}$ and \emph{three} further routes
that consume non-cycle-class data (a Jacobi form, a lattice, a 6d
$\mathcal{N} = (2,0)$ superconformal datum); the Mordell--Weil group
$\mathrm{MW}(\pi)$ of an elliptic fibration $\pi: \mathcal{E} \to \bP^1$
sitting inside $K3 \times E$ indexes the \emph{internal} simple-root
system of the resulting generalised Kac--Moody, \emph{not} an external
orbit symmetry on $H_4(K3 \times E; \ZZ) = \ZZ^{23}$.
\end{corollary}

\subsection{Six-dimensional holomorphic Chern--Simons as the
Costello--Francis--Gwilliam avatar}
\label{wn:subsec:spine-6d-hCS}

\subsubsection{Classical BV datum and the shift law}

\begin{theorem}[Classical BV datum for 6d $\hCS$ on a CY$_3$]
\label{wn:thm:spine-hCS-classical}\ClaimStatusTheorem
Let $X$ be a complex $3$-fold with holomorphic volume form
$\Omega_X \in H^{3,0}(X)$ and let $\fg$ be a finite-dimensional Lie
algebra with non-degenerate symmetric invariant pairing
$\langle\cdot,\cdot\rangle$. The classical BV datum is
\[
\cE_{\hCS} \;=\; \Omega^{0,\bullet}(X, \fg)[1]
\;=\;\Omega^{0,0}[1] \oplus \Omega^{0,1} \oplus \Omega^{0,2}[-1]
\oplus \Omega^{0,3}[-2],\qquad
\cA \;=\; c + A_{0,1} + A^*_{0,2} + c^*_{0,3},
\]
with BRST ghost numbers $1, 0, -1, -2$. The BV pairing
\[
\omega_{\mathrm{BV}}(\alpha, \beta) \;=\; \int_X \Omega_X \wedge
\langle \alpha, \beta\rangle
\]
is of degree $-1$ and non-degenerate via Serre duality between
$\Omega^{0,q}$ and $\Omega^{0,3-q}$. The classical action
\[
S_{\mathrm{cl}}(\cA) \;=\; \int_X \Omega_X \wedge
\bigl\langle \tfrac{1}{2}\cA,\, \bar\partial \cA\bigr\rangle
\;+\; \int_X \Omega_X \wedge
\bigl\langle \tfrac{1}{6}\cA,\, [\cA, \cA]\bigr\rangle
\]
satisfies the classical master equation
$\{S_{\mathrm{cl}}, S_{\mathrm{cl}}\}_{\omega_{\mathrm{BV}}} = 0$.

\emph{Shift law.} On a CY $d$-fold $X$, the BV pairing has degree
$d - 4$: the entries of the canonical table are
\[
\bigl(d,\; \mathrm{shift},\; E_n^{\mathrm{cl}}\bigr) \;\in\;
\bigl\{(2, -2, E_2),\; (3, -1, E_1^{\mathrm{BV}}),\; (4, 0, E_0),\;
(5, +1, E_5\text{-Poisson})\bigr\}.
\]
At $d = 3$ the $(-1)$-shifted symplectic structure quantises to an
$E_3^{\mathrm{hol}}$-factorisation algebra; at $d = 5$ the
$(+1)$-shifted bracket is Poisson, not symplectic
(Pantev--To\"en--Vaqui\'e--Vezzosi $2017$ \S3).
\end{theorem}

\begin{proof}[Proof]
\emph{Well-definedness of $S_{\mathrm{cl}}$.} The real dimension of $X$
is $6$; a top form has bidegree $(3,3)$. For $\cA \in \Omega^{0,\bullet}$
of mixed bidegree, $\Omega_X \wedge \langle \cA, \bar\partial \cA
\rangle$ lands in $(3,3)$ only after summing over all BRST channels
$(0,p) \otimes (0,q) \to (0,p+q+1)$ with $p+q+1 = 3$. The naive
one-form restriction $A \in \Omega^{0,1}$ alone gives bidegree
$(3, 2)$, a $5$-real form, not top --- precisely the reason the
restricted Lagrangian fails in $6$d the way classical Chern--Simons
fails on a $4$-manifold. The BV ghost expansion restores top-degree
integrability.

\emph{Classical master equation.} Quadratic self-pairing:
$\int \Omega_X \wedge \langle \bar\partial\cA, \bar\partial\cA\rangle = 0$
by graded skew-symmetry of $\langle\cdot,\cdot\rangle$ combined with
$\bar\partial^2 = 0$ integrated against $\Omega_X$. Quadratic--cubic
cross term: invariance of the pairing plus $\bar\partial$-Stokes
reduces
$\int \Omega_X \wedge \langle \bar\partial\cA, [\cA, \cA]\rangle$
to a boundary term and cancelling interior pairs. Cubic--cubic
term: Jacobi on $\fg$ forces
$\langle [\cA, \cA], [\cA, \cA]\rangle = 0$ by total antisymmetrisation
on three copies. The three contributions sum to zero; the classical
master equation closes.

\emph{Shift-law derivation.} Serre duality on $\Omega^{0,q}(X, \fg)$
and $\Omega^{0,d-q}(X, \fg)$ pairs to
$\int_X \Omega_X \wedge \langle \cdot, \cdot\rangle$, which has
cohomological degree $d - q - (q) = d - 2q$. Summing over the BRST
expansion shifts total degree by $-4$, giving pairing degree $d - 4$.
At $d = 3$ this is $-1$, $(-1)$-shifted symplectic. At $d = 5$ it is
$+1$, shifted Poisson with bracket of degree $+1$ via
Pantev--To\"en--Vaqui\'e--Vezzosi.
\qedhere
\end{proof}

\subsubsection{Quantum observables and the Bochner--Martinelli
propagator}

\begin{theorem}[BV quantisation of $6$d $\hCS$ on $\CC^3$]
\label{wn:thm:spine-hCS-quantum}\ClaimStatusTheorem
The classical datum of Theorem~\ref{wn:thm:spine-hCS-classical} admits
a perturbative BV quantisation realising a holomorphic factorisation
algebra
\[
\Obs_{\hCS} \colon \mathrm{Open}(\CC^3) \longrightarrow
\mathrm{Ch}_{\mathrm{Dolb}}(\CC),\qquad
\Obs_{\hCS}(\CC^3) \;=\; \bigl(\mathrm{Sym}(\cE_{\hCS}^{\vee}[1])[[\hbar]],\;
Q_{\mathrm{cl}} + \hbar\Delta_{\mathrm{BV}}\bigr).
\]

The free-BV two-point function is
$\langle \cA(z), \cA(w)\rangle_{\mathrm{free}} =
\hbar\, P_{\mathrm{BM}}(z, w)\,\mathrm{Id}_{\fg} \pmod{\bar\partial\text{-exact}}$,
where $P_{\mathrm{BM}}$ is the Bochner--Martinelli propagator
\[
P_{\mathrm{BM}}(z, w) \;=\; \frac{2}{(2\pi i)^3}
\sum_{k=1}^{3} (-1)^{k-1} \frac{\overline{(z_k - w_k)}}{\|z - w\|^{6}}
\,\widehat{d\bar z_k} \wedge dw_1 \wedge dw_2 \wedge dw_3,
\]
with $\widehat{d\bar z_k}$ the $(0,2)$-form omitting the $k$-th factor.
The propagator satisfies $\bar\partial_z P_{\mathrm{BM}} = \delta_\Delta$
in the sense of currents on $\CC^3 \times \CC^3 \setminus \Delta$. The
propagator is the unique $U(3)$-equivariant $t \to 0$ limit of the
heat-kernel propagator $K^{\mathrm{prop}}_t = (\bar\partial^*_{g_0}
\otimes 1) K_t$, up to $\bar\partial$-exact ambiguity; uniqueness comes
from the classification of $U(3)$-equivariant currents on $\CC^3$ with
the correct singular behaviour.

The quantum master equation
$Q_{\mathrm{cl}} S_{\mathrm{eff}}[L] + \hbar\Delta_L S_{\mathrm{eff}}[L]
+ \tfrac{1}{2}\{S_{\mathrm{eff}}[L], S_{\mathrm{eff}}[L]\}_L = 0$
holds at every heat-kernel scale $L > 0$. The renormalisation-group
flow from scale $L$ to $L'$ is a BV automorphism because the
scale-bounded propagator $P_{L'}^L = \int_{L'}^L K_t\, dt$ is itself
$\bar\partial$-exact away from the diagonal, so the flow acts as a
gauge transformation of the BV structure.
\end{theorem}

\subsubsection{The one-loop obstruction: consistent versus covariant,
a Bardeen--Zumino bifurcation}

The one-loop BV anomaly for $6$d $\hCS$ admits two inequivalent
diagrammatic representatives related by a Bardeen--Zumino cochain; both
are correct in their scheme; they must never be conflated.

\begin{theorem}[Consistent and covariant one-loop obstructions]
\label{wn:thm:spine-consistent-covariant}\ClaimStatusTheorem

Let $X$ be a compact CY$_3$ and $\fg$ a semisimple Lie algebra.

\emph{Consistent (BV-cohomology) representative.} The one-loop class
in $H^1_{\mathrm{loc}}(\cE_{\hCS}[-1])$ factorises as
\[
\kanom^{\mathrm{cons}}(X, \fg) \;=\; \hbar\, A(\fg)\,
\frac{\chi_{\mathrm{top}}(X)}{2\,(4\pi)^3}\,\|\Omega_X\|^2_{\mathrm{BCOV}},
\qquad A(\fg) = d^{abc}\, d_{abc}/\dim \fg,
\]
where $A(\fg)$ is the cubic-Casimir coefficient. The class is computed
by the three-leg wheel on $\overline{\mathrm{Conf}}_3(\CC^3)$ with
propagator $P_{\mathrm{BM}}$ on each edge; on a compact $X$ the integral
reduces via Dolbeault--Chern--Weil to
$\int_X c_3(TX) \wedge (\dots) = \chi_{\mathrm{top}}(X)\,
(2(4\pi)^3)^{-1}\, \|\Omega_X\|^2$.

\emph{Covariant (two-leg bubble) representative.} The one-loop
wave-function renormalisation is a distinct graph:
\[
Z^{(1)}_{\cA} \;=\; 1 - \hbar\, C_2(\fg)\,(4\pi)^{-3}
\log(L/\varepsilon),
\]
derived from the two-vertex bubble with quadratic Casimir
$C_2(\fg) = -f^{abc} f_{abc}/\dim\fg$ and logarithmic UV divergence
between scales $\varepsilon < L$.

\emph{Bardeen--Zumino cochain.} There is a holomorphic
Bardeen--Zumino cochain $\mathrm{BZ}^{\mathrm{hol}}(\cA)$ such that
$Q_{\mathrm{BRST}}(\mathrm{BZ}^{\mathrm{hol}})
= \kanom^{\mathrm{cov}} - \kanom^{\mathrm{cons}}$;
the two representatives are cohomologous in
$H^1_{\mathrm{loc}}(\cE_{\hCS}[-1])$ and represent the same obstruction
class.

\emph{The Chern--Simons descent (four-line reconstruction).} The
bubble representative $\kanom^{\mathrm{cons}}$ and the box (three-leg
wheel) representative arise from two different Feynman diagrams
sharing a single BV-cohomology class via the standard Chern--Simons
transgression chain. Step 1: in the classical BV datum, the Chern
character $\mathrm{ch}_3(\cA)^{\mathrm{BV}}
= \tfrac{1}{6}\langle \cA, [\cA, \cA]\rangle$ (cubic Lie bracket)
sits at ghost number $0$ and is $Q_{\mathrm{BRST}}$-closed modulo
$d$-exact terms: this is the local Chern-character
descent $d\mathrm{ch}_3 = 0$, $\mathrm{ch}_3 = d(\mathrm{CS}_3)$
with $\mathrm{CS}_3$ the local Chern--Simons secondary. Step 2: the
box diagram integrates $\mathrm{ch}_3$ against the Bochner--Martinelli
propagator cubed, producing the consistent representative
$\kanom^{\mathrm{cons}} \sim \chi_{\mathrm{top}}(X) A(\fg)$
(topological Euler characteristic times cubic-Casimir coefficient).
Step 3: the bubble diagram integrates the quadratic $\mathrm{CS}_2$
counterterm against the propagator squared, producing a covariant
representative $\kanom^{\mathrm{cov}}$ differing from the consistent
one by $d\mathrm{CS}_2^{\mathrm{bubble}}$. Step 4: the
Bardeen--Zumino cochain $\mathrm{BZ}^{\mathrm{hol}} =
\mathrm{CS}_2^{\mathrm{bubble}} - \mathrm{CS}_3^{\mathrm{box}}$
is precisely the local Chern--Simons difference one level down the
descent tower, and $Q_{\mathrm{BRST}}(\mathrm{BZ}^{\mathrm{hol}})$
equals the difference of the anomaly representatives by direct
application of $Q_{\mathrm{BRST}} \circ \mathrm{CS}_n = d
\mathrm{ch}_n$ on each level.

The reader should note: the descent is not a trick but the standard
Chern--Simons transgression applied one step down to the local
counterterm level. Just as the ordinary Chern--Simons form $\omega_3$
on a $3$-manifold satisfies $d\omega_3 = \mathrm{tr}(F^2)$, and its
variation under a gauge transformation is a $2$-form $\omega_2$
satisfying $\delta\omega_3 = d\omega_2$, the BV descent here is the
functional-integral analog. The bubble and box diagrams live one
level apart on this tower; the BZ cochain is the explicit homotopy
realising the descent on cochains rather than on forms.

\emph{Vanishing tables.} $\kanom^{\mathrm{cons}}$ vanishes
automatically when $\fg \in \{\mathfrak{su}(2), \mathfrak{so}(N), E_6,
E_7, E_8, F_4, G_2\}$ ($d^{abc} = 0$), and hence on every compact
CY$_3$ for these gauge algebras. For $\mathrm{SU}(N \geq 3)$:
$\kanom^{\mathrm{cons}} = 0$ on flat $\CC^3$ (no compact topology), on
$K3 \times E$ (pointwise $c_3(T(K3\times E)) = 0$ via Whitney
$c(TK3)\cdot c(TE)$, hence pointwise--not merely integrated--vanishing),
and on any CY$_3$ with $c_3(TX) = 0$. On the quintic
$Q_5 \subset \bP^4$: $\chi_{\mathrm{top}}(Q_5) = -200$, the anomaly is
nonzero, and trivialisation requires the conjunction of the
Candelas--Horowitz--Strominger--Witten embedding $F_{\cA} = R$ into
$\mathrm{SU}(3)$-tangent holonomy \emph{and} the Green--Schwarz
counterterm $\int B_{0,2} \wedge (\mathrm{tr}\, F^2 - \mathrm{tr}\, R^2)$
sourced by the M-theory $C_3 \wedge X_8$ tadpole inflow. Primary:
Candelas--Horowitz--Strominger--Witten $1985$; Green--Schwarz $1984$;
Costello--Gwilliam--Williams $2021$.

\emph{Higher-loop concentration.} For $n \geq 2$,
$H^1_{\mathrm{loc}}(\cE_{\hCS}[-1])$ at order $\hbar^n$ vanishes
uniformly on semisimple $\fg$ by Whitehead's second lemma
$H^2_{\mathrm{Lie}}(\fg, \fg) = 0$. \emph{Why the lemma applies here.}
Whitehead's lemma states that for every finite-dimensional semisimple
Lie algebra $\fg$ acting on a finite-dimensional module $M$, the
second Lie-algebra cohomology $H^2_{\mathrm{Lie}}(\fg, M) = 0$;
proof idea in one sentence: any $2$-cocycle $\alpha$ satisfies
$d\alpha = 0$, and the Casimir element $C \in U(\fg)$ acts as an
isomorphism on the $\alpha$-defined extension, so $\alpha$ is a
coboundary. The connection to BV anomalies: the order-$\hbar^n$ BV
anomaly obstruction at $n \ge 2$ is classified by
$H^2_{\mathrm{Lie}}(\fg, \fg)$ because the relevant cocycle is the
$\fg$-action on the $\hbar^n$-correction to $S_{\mathrm{eff}}$, which
is an $\fg$-equivariant deformation of the classical action. Whitehead
kills this obstruction for every semisimple $\fg$, so only the
$n = 1$ term ($H^3$-type, not killed by Whitehead) survives as a
genuine obstruction. The one-loop obstruction is the sole BV
anomaly, and its vanishing implies full QME solvability to all
orders in $\hbar$ as a formal power series (convergence is not
claimed).
\end{theorem}

\subsubsection{The $E_3^{\mathrm{hol}}$-structure: holomorphic locality
on the Fulton--MacPherson compactification}

\begin{theorem}[$E_3^{\mathrm{hol}}$ factorisation-algebra structure]
\label{wn:thm:spine-E3-hol-structure}\ClaimStatusTheorem

The quantum observables $\Obs_{\hCS}(\CC^3)$ form a \emph{holomorphic}
factorisation algebra in $\mathrm{Ch}_{\mathrm{Dolb}}(\CC)$, carrying
an $E_3^{\mathrm{hol}}$-algebra structure (Gwilliam--Williams $2021$
\S 2; Costello--Gwilliam $2017$ Vol.~I Thm.~5.3.3 refined to the
holomorphic-locally-constant setting). The topological $E_6$-structure
on $\mathrm{Conf}_2(\CC^3) \simeq_\mathrm{htpy} S^5$ is simply
connected, so the pair-exchange symmetry is trivial on the nose; the
\emph{non-trivial} higher-homotopy coherences sourcing genuine
$E_3^{\mathrm{hol}}$-braiding come from the $\bar\partial$-grading of
the Dolbeault chain complex together with the $\Omega$-background
spectral parameters, not from $\pi_1(S^5) = 0$.

\emph{Correct base space.} The operadic composition structure lives on
the Axelrod--Singer / Fulton--MacPherson compactification
$\FMcpt(n)$ (blow up $\mathrm{Conf}_n(\CC^3)$ along all partial
diagonals), \emph{not} on the naive
$\overline{\mathrm{Conf}}_n(\CC^3)$. On the Fulton--MacPherson base,
the Bochner--Martinelli propagator is absolutely convergent on the
blowup corner charts; on $\overline{\mathrm{Conf}}_n$ convergence is
only conditional. Associativity of the operadic composition is proved
by \v Cech--Dolbeault Mayer--Vietoris on the corner-strata decomposition
of $\FMcpt(n)$ (Axelrod--Singer $1994$; Kontsevich $1999$).

\emph{Chiral Chevalley--Eilenberg avatar.} On objects,
\[
\Obs_{\hCS}(X) \;\simeq\; \Ufact\bigl(\cL^{\hCS}_X\bigr)
\;=\; \CEch\bigl(\cE_{\hCS}, \cO_X\bigr)
\]
where $\cL^{\hCS}_X$ is the local $L_\infty$-space
$\Omega^{0,\bullet}(X, \fg)[1]$ and $\Ufact$ is the
Costello--Gwilliam factorisation envelope functor
(Costello--Gwilliam Vol.~I Ch.~3). On a curve, Francis--Gaitsgory
$2012$ Thm.~1.1 identifies $\Ufact$ with the Beilinson--Drinfeld
chiral Chevalley--Eilenberg; at $d \geq 2$ the BD construction does
not extend directly and $\Ufact$ is the primary definition. The
common notation $\CEch$ denotes the $\Ufact$-construction on Dolbeault
chains in the present monograph.
\end{theorem}

\subsubsection{Minimal $L_\infty$-model: Dolbeault degree-count
formality}

\begin{theorem}[Minimal $L_\infty$-model on flat $\CC^3$]
\label{wn:thm:spine-minimal-model}\ClaimStatusTheorem

The Kontsevich--Soibelman homotopy-transfer minimal model of
$\cE_{\hCS}(\CC^3) = \Omega^{0,\bullet}_c(\CC^3, \fg)[1]$ has
$\ell_n^{\min} = 0$ for all $n \geq 3$. Every homotopy-transfer tree
with at least one internal edge vanishes by a \emph{Dolbeault
degree count}: harmonic inputs sit in $\Omega^{0,0}$, the Hodge
homotopy propagator maps $\Omega^{0,q} \to \Omega^{0, q-1}$, and on
$\Omega^{0, 0}$ its target $\Omega^{0,-1}$ vanishes. This is Kapranov
$1999$ \S 3.1, adapted to flat $\CC^3$. The flat-space claim
"propagator kills the harmonic subspace" is malformed: the
Bochner--Martinelli propagator is a Green operator, not an
annihilation operator; the correct mechanism is the Hodge-homotopy
degree count.

\emph{Formality obstruction on compact CY$_3$.} The cubic bracket is
\[
\ell_3^{\min}(\alpha, \beta, \gamma) \;\sim\;
\int_X \mathrm{At}(TX) \cup \alpha \cup \beta \cup \gamma
\quad\pmod{\text{higher}},
\]
where $\mathrm{At}(TX) \in H^1(X, \Omega^1_X \otimes \mathrm{End}(TX))$
is the Atiyah class of the tangent bundle; the full obstruction is the
Kapranov tower $\{\kappa_n\}_{n \geq 3}$ with leading
$\kappa_3 = \mathrm{At} \cup \mathrm{At}$ and higher Duflo-class Taylor
coefficients (Kapranov $1999$ Thm.~2.8.1; Markarian $2009$).
$\mathrm{At}(TX) = 0$ is a \emph{necessary} formality condition, not
sufficient: the full Kapranov tower must vanish.

\emph{Formality on $K3 \times E$.} The tangent bundle
$T(K3 \times E)$ decomposes as $\pi_1^*TK3 \oplus \pi_2^*TE$;
$\mathrm{At}(TE) = 0$ (elliptic curve is a complex Lie group, tangent
bundle holomorphically trivial) and the K3-direction Kuranishi cubic
receptacle is $H^3(K3, \Lambda^3 T_{K3})$, which vanishes on rank
grounds ($T_{K3}$ has rank $2$, so $\Lambda^3 T_{K3} = 0$). Hence the
full Kapranov tower vanishes and formality holds. The competing
statement "$H^3(K3, \Omega^3_{K3}) = 0$" is an independent rank
vanishing ($\Omega^3_{K3} = 0$ on a surface), not the same vanishing as
$\Lambda^3 T_{K3} = 0$; both are zero but for different reasons, and
the formality argument uses the latter.
\end{theorem}

\subsubsection{Deformation theory and $3$-dualizability}

\begin{theorem}[First-order deformation moduli]
\label{wn:thm:spine-deformations}\ClaimStatusTheorem

The deformation moduli of the $E_3^{\mathrm{hol}}$-algebra
$\Obs_{\hCS}(\CC^3)$ is
\[
\mathrm{Def}(\Obs_{\hCS}) \;=\; \mathrm{HH}^{\bullet}_{E_3}
(\Obs_{\hCS}, \Obs_{\hCS})[3] \;\simeq\; Z_{E_4}(\Obs_{\hCS})
\]
(Francis $2013$ Thm.~2.29, $n$-centre of an $E_n$-algebra).
\emph{The $n$-centre identity, in two sentences.} For an
$E_n$-algebra $A$, the $n$-centre $Z_{E_{n+1}}(A)$ is the universal
$E_{n+1}$-algebra acting centrally on $A$, generalising the classical
fact that the centre of an associative algebra is commutative
($E_1$-algebra has $E_2$-centre). Francis's theorem identifies the
$E_n$-Hochschild cochain complex (which controls $E_n$-deformations)
with the $n$-centre shifted by $n$: $\mathrm{HH}^{\bullet}_{E_n}(A)[n]
\simeq Z_{E_{n+1}}(A)$, so for $n = 3$ the shift is $[3]$ and the
centre lives in the $E_4$-category. The
first-order moduli on $\CC^3$ with simple $\fg$ is
\[
T_0 \cM \;=\; H^{0,3}_{\bar\partial, c}(\CC^3) \otimes
\bigl(\mathrm{Sym}^2(\fg^\vee)^{\fg}\bigr) \;=\;
\CC \cdot (B \otimes \Omega_{\CC^3}),
\]
where $B \in \mathrm{Sym}^2(\fg^\vee)^{\fg}$ is the Killing form.
The moduli match the three-parameter Yangian
$Y_{\epsilon_1, \epsilon_2, \epsilon_3}(\widehat{\fg})$ reduced to the
CY slice $\sum_i \epsilon_i = 0$ by the $(-1)$-shifted symplectic
structure: the BV pairing preservation pins the three Yangian
parameters to a two-parameter family on the CY slice and further to a
one-parameter family via the $S_3$-quotient, matching the
one-dimensional $T_0\cM$ as a structural consequence, not a dimensional
coincidence. The CY-slice relation $\sum\epsilon_i = 0$ lives on the
Yangian side; the single-rank invariant
$\mathrm{Sym}^2(\fg^\vee)^{\fg} = \CC$ is independent of the
$\epsilon$-parameters, forced by Whitehead for simple semisimple $\fg$.
\emph{Why one-dimensional.} For simple semisimple $\fg$, Whitehead's
second lemma (applied to $M = \fg \otimes \fg$) combined with the
existence of a unique-up-to-scalar Killing form $B$ shows
$\mathrm{Sym}^2(\fg^\vee)^{\fg} = \CC \cdot B$: any symmetric
invariant bilinear form on a simple Lie algebra is proportional to
the Killing form, because the action of the Casimir on symmetric
squares has a unique $\fg$-invariant line. This forces the moduli
space of $\fg$-invariant quadratic deformations to be one-dimensional
independent of $\epsilon$-parameters.
\end{theorem}

\begin{theorem}[$E_3$-Koszul self-duality]
\label{wn:thm:spine-E3-koszul}\ClaimStatusTheorem

The $E_3$-operad in chain complexes over $\QQ$ is Koszul self-dual up
to shift:
\[
\mathcal{D}_3^! \;\simeq\; \mathrm{Lie}[2].
\]
\emph{Two-sentence sketch.} The Koszul dual of $E_n$ in chain complexes
is the operad $E_n[n-1]$ shifted by one less than the ambient dimension
(Fresse 2017 Vol.~I Thm.~12.3.A proven via the explicit bar complex on
the little-discs chain operad), and at $n=3$ the shift exposes the
Jacobi identity of a degree-$2$ Lie bracket as the $E_3$-dual of the
Poisson Jacobi identity --- so $\mathcal{D}_3^! \simeq \mathrm{Lie}[2]$
is the self-dual statement with the single Koszul-shift correction.
The reader can verify this at the level of generators: $E_3$ has a
degree-$0$ binary product and a degree-$1$ Poisson bracket, whose
Koszul duals are a degree-$2$ cobracket and a degree-$2$ Lie bracket
respectively, matching $\mathrm{Lie}[2]$.

The strict statement (Gwilliam--Williams $2021$ \S 5.3) and the
homotopy statement (Francis--Gaitsgory $2012$, augmented $E_n$-algebras
in a stable $\infty$-category) are compatible via
Fresse $2017$ Vol.~I Thm.~12.3.A composed with the Positselski
coderived/contraderived transfer (Positselski $2011$ Memoirs AMS
$212$). At $n = 3$ the compatibility requires a three-step chain:
\[
\text{Gwilliam--Williams strict}
\;\xrightarrow{\;\text{Fresse Thm.~12.3.A}\;}
\text{bar--cobar adjunction}
\;\xrightarrow{\;\text{Positselski transfer}\;}
\text{Francis--Gaitsgory}\;\infty\text{-categorical}.
\]
The composition is not a "follows from"; each link is a separate
theorem with its own scope.
\end{theorem}

\begin{theorem}[$3$-dualizability in the abelian sector on flat
$\CC^3$]
\label{wn:thm:spine-3-dual-abelian}\ClaimStatusTheorem

The $E_3^{\mathrm{hol}}$-trace
\[
\mathrm{Tr}^{E_3}_{S^5} : \Obs_{\hCS}(\CC^3) \otimes \Obs_{\hCS}(\CC^3)
\longrightarrow \CC,\qquad
\langle \cO_1(z_1), \cO_2(z_2)\rangle = \int_{S^5}
P_{\mathrm{BM}}(z_1, z_2) \cdot \cO_1 \cO_2,
\]
is nondegenerate on the abelian sector $\fg = \fgl_1$, and
$\Obs_{\hCS}(\CC^3)|_{\fgl_1}$ is $3$-dualisable in
$\mathrm{Alg}_{E_3}(\mathrm{Mod}_{\CC})$ (Lurie \emph{Higher Algebra}
$5.3.2$ applied to the free $E_3$-algebra on the dualizable input
$\Omega^{0,\bullet}_c(\CC^3)[1] \simeq \CC[-3]$).

\emph{Non-abelian failure on $\CC^3$.} For semisimple $\fg$ with
$\dim\fg \geq 3$, Gwilliam--Williams $2021$ Prop.~5.3.2 computes
$\mathrm{HH}^0_{E_3}(\Obs_{\hCS}(\CC^3)|_{\fg}) = \CC[[\tau_1, \tau_2,
\tau_3]]$, infinite-dimensional. The 3-dualizability finiteness
axiom fails. The Calabi--Yau slice $\sum_i \tau_i = 0$ cuts
codimension $1$ to $\CC[[\tau_1, \tau_2, \tau_3]]/(\tau_1 + \tau_2 +
\tau_3)$, still infinite-dimensional; the triality $Y_{\tau_1, \tau_2,
\tau_3}(\widehat{\fgl}_1)$ of the Miki automorphism cyclically
permutes the three formal deformation parameters.
\end{theorem}

\begin{conjecture}[Compact CY$_3$ recovery of 3-dualizability]
\label{wn:conj:spine-compact-recovery}\ClaimStatusConjectured

On a compact CY$_3$ $X$ with a semisimple $\fg$,
$\Obs_{\hCS}(X)|_{\fg}$ is $3$-dualisable in
$\mathrm{Alg}_{E_3}(\mathrm{Mod}_{\CC})$. Equivalently,
$\mathrm{HH}^{\bullet}_{E_3}(\Obs_{\hCS}(X))$ is finite-dimensional in
each cohomological degree. Chain-level evidence on $K3 \times E$ via
Gritsenko $V_1$-lift and Igusa--Borcherds denominator; tree-level
evidence on the quintic; general $(\infty,1)$-categorical proof
requires a compact-$X$ extension of Gwilliam--Williams $2021$
Prop.~5.3.2 not currently in primary literature.
\end{conjecture}

\subsection{Five modular-characteristic subscripts, never conflated}
\label{wn:subsec:spine-five-kappas}

The programme uses five numerical invariants to organise the
CY-to-chiral landscape, and a sixth $\varrho$ anomaly ratio implicit in
the $\mathsf{B}$-row derivation of $K^\kappa = 8$ below. Each
subscript is attached to a specific construction on a specific
object; the subscripts coincide numerically only by accident and only
at special points.

\begin{definition}[Five $\kappa$-subscripts]
\label{wn:def:spine-kappas}\ClaimStatusDefinition
\begin{align*}
\kch(\cA_X) &\;\coloneqq\; \sum_q (-1)^q h^{0,q}(X)
= \chi(\cO_X)\;\;\text{at}\;d\leq 2;\;\;
\text{modified Hodge supertrace at}\;d \geq 3
\;\text{(chiral-side, via }\Phi\text{)};\\
\kcat(X) &\;\coloneqq\; \chi(\cO_X)
\qquad\text{(categorical Euler characteristic, K\"unneth-multiplicative
on products)};\\
\kBKM(\Phi_N) &\;\coloneqq\; \mathrm{wt}(\Phi_N) = c_N(0)/2
\qquad\text{(Borcherds weight of the denominator paramodular form)};\\
\kfib(X, L) &\;\coloneqq\; \mathrm{rk}(L)
\qquad\text{(rank of a fibre or decoration lattice $L$ attached to $X$)};\\
\kanom(X, \fg) &\;\coloneqq\;
\hbar\, A(\fg)\, \frac{\chi_{\mathrm{top}}(X)}{2(4\pi)^3}
\|\Omega_X\|^2_{\mathrm{BCOV}}
\qquad\text{(one-loop BV obstruction of $6$d $\hCS$)}.
\end{align*}
Bare $\kappa$ is forbidden. Subscripts are mandatory at every use; the
subscript identifies which construction is invoked.
\end{definition}

\begin{theorem}[Universal Borcherds-weight identity, at two scopes]
\label{wn:thm:spine-universal-kappa-BKM}\ClaimStatusTheorem

\emph{CHL Borcherds-weight scope.} On the CHL slice
$N \in \{1,2,3,4,6\}$ (those $N$ with $\varphi(N) \mid 2$ admitting a
paramodular witness), the Gritsenko--Nikulin 1998 Thm.~2.1 paramodular
Borcherds lift of the Jacobi form $\phi_{0,1}^{K3, g_N}$ twined by an
order-$N$ K3 symplectic automorphism has constant Fourier coefficient
$c_N(0) \in \{10, 4, 2, 2, 2\}$, and
\[
\kBKM(\Phi_N) \;=\; c_N(0)/2 \;\in\; \{5, 2, 1, 1, 1\}.
\]
\emph{Why weight equals half the constant coefficient.} The Borcherds
multiplicative lift applied to a weight-$0$ index-$1$ Jacobi form
$\phi(\tau, z) = \sum c(n, \ell) q^n \zeta^\ell$ produces an
automorphic form on the signature-$(b^+, 2)$ orthogonal group whose
weight is $c(0, 0)/2$ --- this is Borcherds' 1995 theorem 13.3,
derived by explicit computation: the lift's logarithmic derivative
along the Hodge embedding hits the Eisenstein-like summand whose
normalisation is set by $c(0, 0)$, and the factor of $2$ arises
from the standard $\mathrm{SL}_2 \subset \mathrm{O}(b^+, 2)$
embedding which doubles the weight scale. For the CHL-twined input,
$c_N(0) = c(0, 0)$ of the twined $\phi_{0,1}^{K3, g_N}$; the values
$\{10, 4, 2, 2, 2\}$ come from the Eichler--Zagier computation of
the $M_{24}$-twined K3 elliptic genus constant coefficient, which
equals the $g_N$-character of the weight-$0$ summand.

\emph{Gritsenko additive-lift scope.} On the same CHL slice, the
Gritsenko $1999$ Thm.~1.2 additive lift of a Jacobi form of weight
$k(N) \in \{0, 2, 4, 6, 8\}$ and index $1$ produces a paramodular
form of weight
\[
\mathrm{wt}(\mathrm{Grit}(\phi_{k(N), 1})) \;\in\; \{5, 4, 3, 2, 1\}.
\]

The two ladders $\{5,2,1,1,1\}$ and $\{5,4,3,2,1\}$ are invariants of
\emph{two distinct constructions} on \emph{different Jacobi inputs};
they coincide at $N = 1$ by accident, the numerical match forcing
$\phi_{0,1} = \mathrm{Grit}(\phi_{0,1})$ at that one point. They are
not a single sequence with extensions.

\emph{Non-CHL $N \in \{5, 7, 8\}$ scope.} The rows $N \in \{5, 7, 8\}$
carry half-integer weights $\{1/2, 1/4, 0\}$ that are not Borcherds
weights of $\phi_{0,1}$. They arise from genuine automorphic forms on
the metaplectic cover $\Mpr_4(\mathbb{A})$, classified by
Shimura--Waldspurger lifting through the Weil theta correspondence
$\theta_{2,3}: \Mpr_2 \to \mathrm{O}(3,2)$: a weight-$1$ elliptic cusp
form $g_N \in S_1(\Gamma_0(4N), \chi_N)$ lifts via Shimura $1973$ to
weight $1/2$, then theta-lifts via Waldspurger $1980$ to a paramodular
Siegel form of weight $1/2$ on $\Mpr_4$. The "$1/4$" at $N = 7$ is a
classification index of an order-$4$ \emph{central extension} of
$\Mpr_4$ by $\mu_4$, not a Chern--Simons gerbe class.

\emph{The Lorgat $2020$ identity.} At $N = 1$:
$\frac{1}{64}\Delta_5(2Z) = \Phi(Z)$ as a weight-$5$ paramodular
Borcherds product on $\mathbb{H}_{3,2}$; the divisibility
$64 \mid f(n, \ell, m)$ on the Fourier coefficients of $\phi_{0,1}$ is
\emph{equivalent} to the $\Spr_4(\ZZ)$-equivariance of the
integer-normalised Borcherds product (new structural theorem via
archimedean Eisenstein integrality class in $H^1(\Gamma_1, \ZZ/64)$).
Primary: Borcherds $1995$ Thm.~13.3; Gritsenko--Nikulin $1998$
Thm.~2.1; Borcherds $1998$ Thm.~10.1/13.3.
\end{theorem}

\begin{theorem}[Four-value crystallisation on $K3 \times E$]
\label{wn:thm:spine-four-values}\ClaimStatusTheorem

On $K3 \times E$, the four modular-characteristic values
$\{2, 3, 5, 24\}$ arise from \emph{four distinct constructions}, each
at an explicit scope:
\begin{align*}
\kcat^{\mathrm{K3\text{-}fibre}}(K3 \times E) &\;=\; \chi(\cO_{K3})
\;=\; h^{0,0} + h^{0,2} \;=\; 1 + 1 \;=\; 2
\quad\text{(K3-fibre reading of the Mukai-lattice $\chi$)};\\
\kch(\cH_{\Delta_5}) &\;=\; \mathrm{rk}_{\mathrm{hyp}}(\fg_{\Delta_5})
\;=\; 3
\quad\text{(Cartan-rank of the rank-$3$ hyperbolic BKM $\fg_{\Delta_5}$ on
$\mathrm{II}_{3,2}$)};\\
\kBKM(\Phi_1) &\;=\; c_1(0)/2 \;=\; 5
\quad\text{(Borcherds weight of $\Delta_5$, via Gritsenko--Nikulin
$1998$)};\\
\kfib(K3, \Wttl(K3)) &\;=\; \mathrm{rk}(\Wttl(K3)) \;=\; 24
\quad\text{(Mukai-lattice rank, signature $(4, 20)$)}.
\end{align*}

\emph{Scope declarations, and the two parallel readings used elsewhere.}
Each value above is attached to a specific construction; two natural
alternative readings arise for $\kcat$ and $\kch$ and must be
disambiguated whenever invoked:

\begin{center}
\begin{tabular}{lll}
\toprule
Subscript & K3-fibre scope & Total-space scope \\
\midrule
$\kcat^{\mathrm{K3\text{-}fibre}}$ & $2 = \chi(\cO_{K3})$ &
$\kcat^{\mathrm{total}}(K3\times E) = 0 = \chi(\cO_{K3})\cdot\chi(\cO_E) = 2\cdot 0$ \\
\multicolumn{3}{l}{\emph{Both correct; K\"unneth-multiplicative
on products; the $K3$-fibre reading is the manifesto default.}} \\
\midrule
Subscript & Cartan-rank reading & Mukai-enhanced Heisenberg reading \\
\midrule
$\kch$ & $3 = \mathrm{rk}_{\mathrm{hyp}}(\fg_{\Delta_5})$ &
$\kch^{\mathsf{B}\text{-row}}(\cH_{\Mukr}(K3)) = 4 = c_+(\Mukr(K3))$ \\
\multicolumn{3}{l}{\emph{Both correct; Cartan-rank reading tracks the
$\fg_{\Delta_5}$ BKM; Mukai-enhanced reading tracks the $\mathsf{B}$-row}} \\
\multicolumn{3}{l}{\emph{of the five-archetype ceiling via
$K^{\kch}_{\mathsf{B}} = \varrho K = (1/6)\cdot 48 = 8 = 2c_+$.}} \\
\bottomrule
\end{tabular}
\end{center}

\emph{The one-loop anomaly vanishes pointwise.} By Whitney on the
split tangent bundle, $c(T(K3 \times E)) = c(TK3) \cdot c(TE) =
(1 + c_2(TK3)) \cdot 1 = 1 + c_2(TK3)$, so $c_3(T(K3 \times E)) = 0$
as a cohomology class; $\kanom(K3 \times E, \fg) = 0$ for every $\fg$,
pointwise on the integrand, not merely after integration.
\end{theorem}

\subsection{The Igusa $\Delta_5$ generalised Kac--Moody
$\fg_{\Delta_5}$}
\label{wn:subsec:spine-delta5}

\begin{theorem}[Structure of $\fg_{\Delta_5}$]
\label{wn:thm:spine-gDelta5}\ClaimStatusTheorem

$\fg_{\Delta_5}$ is a Borcherds--Kac--Moody Lie superalgebra with
denominator identity
\[
\Phi(Z) \;=\; \sum_{w \in W^{(2)}} \det(w)\,
e^{-2\pi i (w\rho, Z)}
\biggl(1 - \sum_a m(a)\, e^{-2\pi i\, (w(\rho + a), Z)}\biggr),
\qquad
Z \in \mathbb{H}_{3,2},
\]
and the following root data:
\begin{itemize}
\item \emph{Real simple roots} $\{\delta_1, \delta_2, \delta_3\}$ with
Gram matrix $\mathrm{diag}(2,2,2) - 2(E - I)$, rank-$3$ hyperbolic,
eigenvalues $\{+4, +4, -2\}$. Weyl vector
$\rho = \tfrac{1}{2}(\delta_1 + \delta_2 + \delta_3)
= f_2 - \tfrac{1}{2}f_3 + f_{-2}$.
\item \emph{Even imaginary simple roots}
$\Delta^{\mathrm{im}}_0 = \{\tau(a) \cdot a :
(a, a) = 0,\; \tau(a) > 0\}$ with $\tau(a_0) = 9$.
\item \emph{Odd imaginary simple roots}
$\Delta^{\mathrm{im}}_1 = \{m(a) \cdot a :
(a, a) < 0,\; m(a) < 0\}$
with $m(a) = -\tfrac{1}{64} f(n, \ell, m)$.
\item \emph{Root multiplicities} $\mathrm{mult}_\alpha = f(nm, \ell)$,
Kac asymptotic $f(n, \ell) = O(\exp \sqrt{4n - \ell^2})$.
\item \emph{Real-root subalgebra} is the Feingold--Frenkel $F_3$
rank-$3$ hyperbolic Kac--Moody (\emph{not} $\widehat{\mathfrak{sl}}_3$).
$\fg_{\Delta_5}$ is the super-completion of $F_3$ by Borcherds
super-Serre adjunction of odd imaginary simple roots. \emph{Why $F_3$
and not affine $\widehat{\mathfrak{sl}}_3$.} $F_3$ has Cartan of
signature $(2, 1)$ (two positive, one negative), which is the
\emph{hyperbolic} type in Kac's classification; $\widehat{\mathfrak{sl}}_3$
has Cartan of signature $(2, 0, 1)$ with a null direction (the
centre), placing it in the \emph{affine} class. The eigenvalues
$\{+4, +4, -2\}$ of the $\Delta_5$ Gram matrix show two positive and
one negative, matching $F_3$'s hyperbolic signature and ruling out
the affine possibility. The reader can distinguish the two classes
by checking whether the Cartan bilinear form has a null vector:
affine algebras do (the imaginary root $\delta$), hyperbolic
algebras do not.
\item \emph{Group embedding.}
$\Spr_4(\ZZ)/\{\pm I\} \hookrightarrow \SOr(\Lambda^{3,2})$ as a
finite-index subgroup (\emph{not} an isomorphism: the exterior-square
$\wedge^2$ lands in $\mathrm{SO}$ of a rank-$6$ lattice of signature
$(3,3)$, descending to $(3,2)$ via Pl\"ucker $\mathrm{LG}(2,4)
\hookrightarrow \bP(\wedge^2 V_4)$, and the integral refinement has
finite cokernel).
\end{itemize}
\end{theorem}

\begin{theorem}[Dimension-stratified Borcherds census]
\label{wn:thm:spine-dimension-census}\ClaimStatusTheorem

The three BKM algebras organising the Vol~III landscape live on
\emph{three distinct lattices}, at \emph{three distinct CY dimensions},
with \emph{three distinct Borcherds lifts}. They are unified by the
universal weight identity $\kBKM(\Phi) = c(0)/2$ and by a
Borcherds-functorial denominator restriction, \emph{not} by synthetic
envelope enclosure.

\begin{center}
\begin{tabular}{lllll}
\toprule
BKM & Lattice & Cartan rank & CY dim & Denominator $\Phi$ \\
\midrule
$\fg_{\Delta_5}$ & $\Lambda^{3,2} = \mathrm{II}_{3,2}$ & $3$ &
$d = 3$ on $K3 \times E$ & $\Delta_5$ weight $5$ \\
Borcherds Monster $\fm$ & $\mathrm{II}_{1,1}$ (virtual) & $26$ &
$d = 3$ virtual CY & $j(p) - j(q)$ weight $0$ \\
Fake Monster $\fg_{\mathrm{FM}}$ & $\mathrm{II}_{25,1}$ & $26$ &
$d = 5$ on $K3 \times K3 \times E$ & Borcherds $\Phi_{12}$ weight $12$\\
\bottomrule
\end{tabular}
\end{center}

At $d = 3$ on $K3 \times E$, the rank count
$\mathrm{rk}(\LLeech) = 24 > h^{1,1}(K3) = 20$ obstructs a
Leech-lattice-supported BKM, forcing the Fake Monster to $d = 5$ where
$h^{2,2}(K3 \times K3) = 404$ is plentiful (direct Hodge-diamond:
$h^{0,0}\cdot h^{2,2} + h^{2,2} \cdot h^{0,0} + h^{1,1}\cdot h^{1,1}
+ h^{2,0}\cdot h^{0,2} + h^{0,2}\cdot h^{2,0}
= 1 + 1 + 400 + 1 + 1 = 404$).

The Borcherds Monster $\fm$ lives on $\mathrm{II}_{1,1}$, outside the
Borcherds $1998$ Grassmannian singular-theta lift (which requires the
signature-$(b^+, 2)$ Hermitian symmetric case $b^+ \geq 3$). Frenkel--
Lepowsky--Meurman $1988$ orbifold the Leech lattice VOA
$V_{\LLeech}$ at a $\ZZ/2$-involution to produce the Moonshine module
$V^\natural$; this is an orbifold of a \emph{lattice vertex algebra},
not of a Calabi--Yau variety. Consequently the notation
``$V^\natural = \mathrm{Sp}_{T^{24}_{\LLeech}/\ZZ_2, E_\tau}
(\Phi^{\mathrm{FA}}_3(X_B))$'' of an earlier draft is incorrect.
\emph{The dimensional check the reader should perform.} For
$\mathrm{Sp}^{\mathrm{ch}}_{\Sigma_{d-1}, C}$ to apply at $d = 3$,
we need the CY host to be three-complex-dimensional; but
$T^{24}_{\LLeech}/\ZZ_2$ is a $24$-complex-dimensional orbifold (it is
the Leech torus, a $24$-dimensional complex torus, modulo a
$\ZZ/2$-involution that does not reduce complex dimension), so
$(T^{24}_{\LLeech} \times E)/\ZZ_2$ is $25$-complex-dimensional, not
$3$. The Stage-$2$ specialisation functor cannot produce a curve-level
chiral algebra from a $25$-fold host. The correct statement is that
$V^\natural$ is the chiral-boundary shadow of a virtual CY-$3$ datum
(the Monster lives on $\mathrm{II}_{1,1}$, and the ``virtual CY-$3$''
means we treat $\mathrm{II}_{1,1}$ formally as if it arose from a
three-fold, tracking only the automorphic-denominator data), not a
genuine Stage-$2$ specialisation of $\Phi^{\mathrm{FA}}_3$ of a
geometric variety.

\emph{Hyperbolic restriction correspondence.} The genuine Borcherds
$1998$ \S 14 K\"unneth input identifies $\Phi_{12}$ on
$\mathrm{II}_{26,2}$ restricted along
$\mathrm{II}_{2,2} \hookrightarrow \mathrm{II}_{25,1}$ with $\Phi_{10}$
on $\mathrm{II}_{3,2}$, and $\Phi_{10} = \Delta_5^2$. This is the
arithmetic content behind the naive "envelope" statement, and it lives
entirely at $b^- = 2$; the rank-$(3,3)$ "envelope" at $\Lambda^{3,3}$
is not the locus where a holomorphic automorphic form exists.

\emph{Hauptmodul property.} The Hauptmodul property of the
McKay--Thompson series is Borcherds $1992$ (Hecke-replication on
$\fm$'s root lattice), \emph{not} a consequence of the
$\Phi^{\mathrm{FA}}_3$ Stage-$1$ contractibility. The two are
categorically independent arithmetic inputs.
\end{theorem}

\subsection{CoHA, Miki triality, and the chiral Yangian}
\label{wn:subsec:spine-CoHA-Miki}

\begin{theorem}[$\CoHA(\CC^3)$, three presentations, Miki $S_3$-triality]
\label{wn:thm:spine-coha-miki}\ClaimStatusTheorem

\emph{Three presentations of the same algebra, after shuffle-localisation.}
Over the localised ring $\mathbb{F} = \CC(\epsilon_1, \epsilon_2)$ with
$\epsilon_3 = -\epsilon_1 - \epsilon_2$, the three presentations of the
positive half of the affine Yangian agree:
\[
\underbrace{\CoHA(\CC^3) \otimes_{\CC[\epsilon_1,
\epsilon_2]} \mathbb{F}}_{\text{Schiffmann--Vasserot shuffle}}
\;\simeq\;
\underbrace{Y^+_{\epsilon_1, \epsilon_2, \epsilon_3}
(\widehat{\fgl}_1)}_{\text{Drinfeld currents, Tsymbaliuk 2017}}
\;\simeq\;
\underbrace{\mathcal{SH}^{+, \mathrm{MO}}(\CC^3)}_{\text{Maulik--Okounkov
RLL / Negu\c t shuffle}}.
\]
\emph{Why three presentations, and how they match.} Each presentation
names the same algebra through a different generator basis and
relation set: (a) the Schiffmann--Vasserot shuffle presents $\CoHA$
via the shuffle-product kernel $\omega(z, w)
= \prod_i (z - w - \epsilon_i)/(z - w)^3$ acting on polynomials; the
algebra is the localised polynomial ring with shuffle multiplication.
(b) The Drinfeld currents presentation writes $Y^+$ generators as
$e_i(z), h_i(z), f_i(z)$ with rational function relations; Tsymbaliuk
2017 provides an explicit isomorphism
$\Psi_{\mathrm{SV-Dr}}\colon \CoHA \otimes \mathbb{F} \to Y^+$
matching the two generator sets by comparing their action on the
Nakajima Fock module. (c) The Maulik--Okounkov RLL presentation builds
$\mathcal{SH}^+$ from the equivariant cohomology of Hilbert schemes of
points on $\CC^2$ with shuffle structure induced from the
Maulik--Okounkov $R$-matrix; the isomorphism with $Y^+$ is Negu\c t's
theorem 2015. The localisation is needed because the isomorphisms
use inverse factors like $(\epsilon_1 \epsilon_2 \epsilon_3)^{-1}$ to
clear denominators on one side of the generator maps. At the
unlocalised integer level, $\CoHA \hookrightarrow Y^+$ with
cokernel annihilated by $\epsilon_1 \epsilon_2 \epsilon_3$ because
the inverse maps acquire these three factors as denominators.

The integer-unlocalised $\CoHA(\CC^3)$ injects into
$Y^+$ over $\CC[\epsilon_1, \epsilon_2]$ with cokernel annihilated by
$\epsilon_1 \epsilon_2 \epsilon_3$. The strict identity
"$\CoHA(\CC^3) = Y^+$" is therefore a localised-isomorphism
statement.

\emph{Miki $S_3$-triality.} The shuffle kernel
$\omega(z, w) = \prod_{i = 1}^{3}(z - w - \epsilon_i)/(z - w)^3$ is
manifestly $S_3$-symmetric under permutations of $(\epsilon_1,
\epsilon_2, \epsilon_3)$. This single symmetry descends to:
\begin{itemize}
\item a Hopf-algebra automorphism of the Drinfeld double
$Y = Y^+ \otimes Y^0 \otimes Y^-$ (Miki $2007$);
\item a shuffle-product automorphism of the positive half $Y^+$
(Schiffmann--Vasserot $2013$ via Tsymbaliuk $2017$);
\item an image automorphism of $\Winf[\lambda]$ under the
evaluation morphism $\mathrm{ev}_\lambda$ (Feigin--Jimbo--Miwa--Mukhin
$2016$).
\end{itemize}
$\Winf[\lambda]$-triality is the $\mathrm{ev}_\lambda$-image shadow of
$Y^+$-triality, not its source. At the Calabi--Yau slice
$\sum \epsilon_i = 0$ (codimension-$1$ hyperplane in the three-parameter
$\epsilon$-space) the triality survives faithfully; the recursion does
\emph{not} truncate, by the Chebyshev relation
$\mathsf{X}_{n+1}(z) = \mathsf{X}_1(z)\mathsf{X}_n(z - \epsilon_3)
- \mathfrak{q}\mathsf{X}_{n-1}(z - 2\epsilon_3)$
with discriminant $\mathsf{X}_1^2 - 4\mathfrak{q}$ nowhere vanishing
on the CY slice.

\emph{Ran-level avatar.} Descent of the $S_3$-triality to a
factorisation algebra on $\mathrm{Ran}(\CC)$ is \emph{conjectural}
(Beilinson--Drinfeld Ch.~3.4 provides the vertex-algebra-to-factorisation-algebra
equivalence; the $S_3$-compatibility of the three projections
$\pi_i : \mathrm{Ran}(\CC) \to \mathrm{Ran}(\CC)$ attached to the three
shuffle directions is an open structural statement).
\end{theorem}

\begin{theorem}[$5$D $\hCS$ to affine Yangian, at honest scope]
\label{wn:thm:spine-5d-hCS-yangian}\ClaimStatusTheorem

\emph{Proved, finite Yangian, perturbatively.} For every simply-laced
semisimple $\fg$ (ADE), the perturbative boundary of $5$D $\hCS(\fg)$
is isomorphic to the $\hbar$-formal finite Yangian
$Y_\hbar(\fg) \otimes \mathrm{Heis}$. \emph{Six-step proof, explicit.}
(1) Costello 2013 quantises $5$D hCS on $\CC^2 \times \RR$ to an
$E_2^{\mathrm{hol}}$-algebra of observables, with the Bochner--Martinelli
propagator on $\CC^2$ providing explicit Feynman amplitudes;
(2) the boundary $\CC^2 \times \{0\}$ carries induced
$E_2^{\mathrm{hol}}$-observables whose current OPE matches the
Maulik--Okounkov $R$-matrix at the Yangian level; (3) Francis 2013
Thm.~2.29 identifies $E_2$-Hochschild cohomology with the
$E_3$-centre, letting us pass from boundary observables to the
universal enveloping Yangian; (4) Whitehead's second lemma
(as reconstructed in the Higher-loop concentration above) closes
the perturbative expansion: no obstructions arise beyond one loop for
semisimple $\fg$; (5) Chan--Paton minuscule embedding realises
$A_n, D_n, E_6, E_7$ Yangians via a fundamental minuscule
representation --- the Yangian generators are built from the
Chan--Paton factors coupled to the $5$D hCS gauge fields on the
boundary; (6) for $E_8$ the embedding fails (no minuscule rep of
$E_8$), replaced by the Kostant--Slodowy slice construction, which
produces $\cW_{\mathrm{fin}}(\fe_8)$ as the Yangian image at the
principal nilpotent orbit rather than the full $Y(\widehat{\fe}_8)$;
the strict inclusion $\cW_{\mathrm{fin}}(\fe_8) \subsetneq
Y(\widehat{\fe}_8)$ is the reader's diagnostic that $E_8$ is the
only exceptional case requiring the slice replacement. ``All orders
in $\hbar$'' is formal, not convergent.

\emph{Proved, abelian affine, Heisenberg fibre.} The abelian
$\fg = \fgl_1$ all-orders affine Yangian
$Y^+_{\epsilon_1, \epsilon_2, \epsilon_3}(\widehat{\fgl}_1)$
Heisenberg sector is produced via Costello 2013 + Costello-Gaiotto
$2018$ + Prochazka-Rapcak $2018$ truncation
$\cW_\infty[\lambda] \twoheadrightarrow \cW_n$.

\emph{Conjectural, non-abelian affine.} The extension to
$Y(\widehat{\fg})$ at non-abelian semisimple $\fg$ is \emph{conjectural};
the obstruction is that the Drinfeld double requires geometric
$\CC^3 \sqcup (\CC^3)^{\mathrm{op}}$ along $S^5$, not merely a
half-space.

\emph{Corner at finite $n$.} The boundary corner $Y_{0,0,n}$ relates to
$Y_{n,n,n}$ not via an $n \to \infty$ inductive limit, but via a
principal Miura quotient (Prochazka--Rapcak $2018$ Thm.~4.5). The earlier
corner-agreement statement inverted the arrow direction.

\emph{Chiral-Yangian disambiguation.} The chiral Yangian
$\mathcal{Y}^{\mathrm{ch}}(\fg)$ is the $L_0$-bounded sub-VOA of the
corner $\mathcal{Y}_{0,0,n}[\psi] = \cW_n[\psi]$, inheriting central
charge $c_n[\psi] = (n-1)(1 - n(n+1)(\psi + \psi^{-1} - 2))$. It is a
\emph{sub-VOA}, not a quotient, not a module.
\end{theorem}

\subsection{The Heisenberg--Mukai $\eta^{-48}$ identity on $K3\times E$}
\label{wn:subsec:spine-heis-mukai}

\begin{theorem}[Heisenberg--Mukai identity, all orders]
\label{wn:thm:spine-heis-mukai}\ClaimStatusTheorem

\[
\chi_{\mathrm{Heis}(\Wttl(K3)^{\oplus 2})}(q) \;=\;
\prod_{n \geq 1}(1 - q^n)^{-48}
\;=\; \frac{1}{\eta(q)^{48} \cdot q^{-2}},
\]
\emph{Direct derivation (three steps a reader can reproduce).}
Step 1: the Heisenberg chiral algebra on a rank-$r$ even lattice $L$
has character $\prod_{n \ge 1}(1 - q^n)^{-r}$: the character counts
bosonic Fock-space states, and each of the $r$ independent lattice
directions contributes an independent bosonic tower generating
$1/\prod(1 - q^n)$. For the rank-$48$ lattice $\Wttl(K3)^{\oplus 2}$,
$r = 48$, giving $1/\eta^{48}$ up to the $q^{-r/24} = q^{-2}$
modular-anomaly prefactor from the normalisation
$\eta(q) = q^{1/24}\prod_{n\ge 1}(1 - q^n)$.
Step 2: no null vectors appear because the Shapovalov form on the
Heisenberg Fock space is built block-diagonally from the lattice
Gram matrix, which for $\Wttl(K3)^{\oplus 2}$ is non-degenerate
(Mukai pairing is even unimodular of signature $(4, 20)$; doubling
preserves non-degeneracy). Step 3: the character is therefore
$1/\eta^{48}$ at all orders in $q$, without minimal-model truncations.
Explicit low-order coefficients
$(g_n)_{n \geq 0} = (1, 48, 1224, 21952, 309876, \dots)$ follow by
direct Euler expansion $\prod_{n\ge 1}(1 - q^n)^{-48}
= \sum_{n \ge 0} g_n q^n$; the reader can check $g_1 = 48$ (the
$48$-dimensional weight-$1$ lattice) and $g_2 = 48 \cdot 49/2 + 48
= 1176 + 48 = 1224$ (symmetric square plus the $L_{-2}$ states).
The explicit value at $g_{24}$ is recomputed on demand by direct
Euler expansion and is not re-inscribed here.

\emph{Match with $\Delta_5^{-2}$ at the Humbert divisor.} Using
Gritsenko--Hulek $1997$ expansion
$\Delta_5(\tau, z, \sigma) = (2\pi i z)\,\eta(\tau)^{12}\eta(\sigma)^{12}
+ O(z^3)$,
\[
\Bigl[(2\pi i z)^2\, \Delta_5^{-2}\Bigr]\Big|_{z = 0}
\;=\; \frac{1}{\eta(\tau)^{24}\,\eta(\sigma)^{24}},
\]
and on the diagonal $\tau = \sigma$ this equals $1/\eta(q)^{48}$. The
$(2\pi i z)^2$ prefactor absorbs the $z \to 0$ collapse along the
Humbert divisor $H_1 \subset \cA_2$; the residue is a two-variable
function of $(\tau, \sigma)$, and single-variable collapse requires
the diagonal specialisation.

\emph{Retraction of the $\cV_{24} = H^0_{\mathrm{DS}}
(L_{-2 + 1/22}(\fsl_2)^{\otimes 22})$ route.} Direct Drinfeld--Sokolov
computation:
$c_{\mathrm{DS}}(L_{-2 + 1/22}(\fsl_2)) = 1 - 6\cdot 21^2/(22\cdot 1)
= 1 - 6(21/22)^2/(1/22)\cdot\mathrm{fractional}\ldots
= -1312/11$ per copy; $22$ copies give $c = -2624$, not $-214$. The
"arithmetic error $-14432/121 \to -1312/11$" alluded to in earlier
revisions is illusory: $-14432/121 = -1312\cdot 11/121 = -1312/11$. The
claim that $\cV_{24}$ arises as an iterated Drinfeld--Sokolov reduction
is retracted. The correct lane is the \emph{positive-definite} lattice
VOA $V_{\Mukr(K3)^{\oplus 2}}$ at central charge $c = 48$; the direct
$\eta^{-48}$ identity bypasses any Virasoro $(2, 45)$-minimal apparatus.
\end{theorem}

\subsection{AdS$_3$ throat and dyon degeneracies}
\label{wn:subsec:spine-AdS3}

\begin{theorem}[Near-horizon throat and dyon count]
\label{wn:thm:spine-AdS3}\ClaimStatusTheorem

\emph{Geometry.} The near-horizon throat of the IIB
D$1$-D$5$-KK-momentum system on $K3 \times S^1 \times \widetilde{S^1}$
at the CHL orbifold point is
$\mathrm{AdS}_3 \times S^3/\ZZ_N \times K3^{g_N}$, with $K3^{g_N}$
the $g_N$-orbifold of K3 by an order-$N$ symplectic automorphism.

\emph{Central charge.} The reduced Brown--Henneaux central charge of
the boundary chiral algebra is $c_L^{\mathrm{reduced}} = 24$ uniformly
in $N$. The quantity $k_N = 24/(N+1) - 2$ is the \emph{Borcherds weight}
of the CHL orbifold's paramodular lift $\Phi_{k_N}$, not a central
charge; the earlier formula $c_N = 24 k_N$ is wrong (it would give
$c_1 = 24 \cdot 5 = 120$ or $240$, contradicting the MSW microstate
relation $c_L = 6k + 24$ with $k = N + 1$).

\emph{Dyon count.} The $\tfrac{1}{4}$-BPS dyon degeneracy at CHL point
$\ZZ_N$ is
\[
d_N(Q, P) \;=\; \oint_{\mathcal{C}_N}
\frac{e^{-i\pi (Q, \Omega) \cdot T (Q, \Omega)^T}}{\Phi_{k_N}(\Omega)^2}\,
d\rho\, d\sigma\, dv,
\]
with leading entropy
$\log d_N = 2\pi\sqrt{\Delta/N} + (k_N + 2)\log\Delta + \cdots$
(Sen $2008$; Dijkgraaf--Verlinde--Verlinde $1997$;
Shih--Strominger--Yin $2005$).
\emph{Why the Siegel contour integral gives this asymptotic.} The
integrand $1/\Phi_{k_N}^2$ is a meromorphic Siegel modular form whose
leading pole at $v = 0$ dominates the contour integral in the
saddle-point limit; saddle-point gives
$\log d_N \sim 2\pi\sqrt{\text{norm}}$ where the norm is
$(Q, P)$-discriminant $\Delta$ rescaled by the CHL order $N$ (the
rescaling comes from the $\Gamma^{(2)}_N$-invariant inner product on
the paramodular cover, where the $N$-fold cover introduces a $1/N$
factor in the quadratic form normalisation). The subleading
$(k_N + 2)\log\Delta$ term is the logarithmic correction from the
$\Phi_{k_N}^2$ weight: a weight-$k$ cusp form contributes
$k \log \Delta$ to the asymptotic expansion; the ``$+ 2$'' is the
standard volume-Jacobian correction in genus-$2$ Siegel-upper-half
space.

\emph{Graviton finiteness.} The chiral-boundary algebraic rigidity
matches the $\mathrm{HH}^2_{E_2}$-vanishing of the $E_2$-chiral algebra
$\Phi_2^{\mathrm{FA}}(D^b\mathrm{Coh}(K3))$ at the Mukai-enhanced
Heisenberg sector, giving a single structural witness
\emph{graviton finiteness $\equiv$ $E_2$-chiral rigidity}.
\end{theorem}

\subsection{Derived-centre complementarity: the five-archetype
landmark ceiling}
\label{wn:subsec:spine-complementarity}

\begin{theorem}[Five-archetype landmark ceiling, scope declared]
\label{wn:thm:spine-five-archetype}\ClaimStatusTheorem

On the canonical seven-witness landmark table
$\{\cH_k,\; \widehat{\fg}_k,\; \beta\gamma_\lambda,\; \mathrm{Vir}_c,\;
\cW_3^k,\; \mathrm{BP}_k,\; \cH_{\Mukr}(K3)\}$, the bridge identity
\[
K^{\kch}(A) \;=\; \varrho(A) \cdot K(A),\qquad
\varrho(A) := \kch(A)/c(A),\;\;
K(A) := c(A) + c(A^!),\;\;
K^{\kch}(A) := \kch(A) + \kch(A^!),
\]
takes values
$K^{\kch} \in \{0, 8, 13, 250/3, 98/3\}$.

This is a \emph{landmark ceiling}, not a universal bound: principal
$\cW_N^k$ at $N \geq 4$ produces further rational values outside the
five-element bucket (Theorem derived-centre-complementarity-strengthened,
Vol I chiral\_center\_theorem.tex eq.~thm-C-full-set). Scope
declaration is mandatory.

\emph{$\mathsf{B}$-row witness (Mukai-enhanced K3 Heisenberg).} The
rank-$24$ Heisenberg chiral algebra on the Mukai lattice
$\Wttl(K3) \simeq E_8(-1)^{\oplus 2} \oplus U^{\oplus 3}$ of signature
$(4, 20)$ has $\varrho = 1/6$ (Mukai-enhanced anomaly ratio, identical
to the Bershadsky--Polyakov minimal-nilpotent DS class), $K = 48 =
2 \cdot \mathrm{rk}(\Wttl(K3))$, hence
\[
K^{\kch}(\cH_{\Mukr}(K3)) \;=\; (1/6) \cdot 48 \;=\; 8
\;=\; 2\, c_+(\Mukr(K3)).
\]

\emph{Three faces of $8$, unified.}
\begin{enumerate}[label=(\roman*)]
\item \emph{Mukai face, categorical.}
$8 = 2c_+(\Mukr(K3))$ via the Beilinson--Drinfeld Koszul-conductor
identity on the $(4, 20)$-signature.
\item \emph{Humbert-monodromy face, arithmetic.}
$8 = \mathrm{ord}(\mathrm{monodromy}\, \cL^{\Delta_5}|_{H_1})$ via
Bruinier $2002$ Prop.~5.1 Heegner-Chern-class reciprocity: the
paramodular multiplier order is $2$, the $\eta^9\vartheta_1$ Jacobi
weight $9/2$ contributes index $4$; total is $2 \cdot 4 = 8$.
\item \emph{Lusztig face, quantum-group.}
$8 = \ell$ where $\zeta^8 = 1$ is the root-of-unity order at which the
Hall--Drinfeld double specialises to the small quantum group
$\mathfrak{u}_\zeta$ (Lusztig $1990$, Drinfeld $1990$ quasi-Hopf scaling
$\hbar = 2\pi i/\ell$).
\end{enumerate}
Faces (i) $\leftrightarrow$ (ii) via Bruinier reciprocity; (ii)
$\leftrightarrow$ (iii) via Drinfeld scaling; (i) $\leftrightarrow$
(iii) follows.

\emph{Independence from $\kBKM$.} The Borcherds weight $\kBKM =
5 = c_1(0)/2$ is a fourth, disconnected invariant of $\Delta_5$ that
lives on the \emph{automorphic-form lane}; it does not enter the
bridge-identity sum. The earlier reading
$K^\kappa = \kch + \kBKM = 3 + 5$ confused two lanes: the
Verdier-dual invariant $\kch^!$ and the Borcherds weight $\kBKM$ are
different quantities; their numerical near-match at $N = 1$ (both $5$)
is accidental.
\end{theorem}

\subsection{Retractions with true hidden structure}
\label{wn:subsec:spine-retractions}

\begin{theorem}[The twenty-agent retraction ledger]
\label{wn:thm:spine-retractions}

Each wrong claim is stated, the precise error identified, and the
ghost-theorem extracted.

\begin{enumerate}
\item \textbf{$\widehat{\mathfrak{sl}}_3 \hookrightarrow \fg_{\Delta_5}$
from $\eta^9$.}
\emph{Error.} Cartan mismatch (rank-$2$ positive-definite vs rank-$3$
hyperbolic); denominator mismatch ($\eta^{11}$ for
$\widehat{\mathfrak{sl}}_3$ vs $\eta^9$ in $\Delta_5$).
\emph{Ghost.} Real-root subalgebra of $\fg_{\Delta_5}$ is the
Feingold--Frenkel $F_3$ rank-$3$ hyperbolic Kac--Moody;
$\eta^9 = \eta^1 \cdot \eta^8$, not $\eta^6 \cdot \eta^3$ (Feingold--
Frenkel $1983$).

\item \textbf{$\mathrm{VOA}[T_{24}] = L_{-6}(\mathfrak{e}_8)$.}
\emph{Error.} Central-charge mismatch $c = -62 \neq -214$.
\emph{Ghost.} The correct avatar of $\cV_{24}$ is the
\emph{positive-definite} lattice VOA $V_{\Mukr(K3)^{\oplus 2}}$ at
$c = 48$; the iterated Drinfeld--Sokolov presentation
$H^0_{\mathrm{DS}}(L_{-2 + 1/22}(\fsl_2)^{\otimes 22})$ is itself
retracted (see Theorem~\ref{wn:thm:spine-heis-mukai}).

\item \textbf{$\kBKM = \kch + \chi(\cO_{\mathrm{fiber}})$ universally.}
\emph{Error.} Accidentally holds at $N = 1$; fails at $N \geq 2$
(at $N = 2$: LHS $= 4$, RHS $= \kch(K3\times E) + \chi(\cO_E) =
0 + 0 = 0$).
\emph{Ghost.} Universal Borcherds-weight identity
$\kBKM(\Phi_N) = c_N(0)/2$ (Theorem~\ref{wn:thm:spine-universal-kappa-BKM}).

\item \textbf{Fake Monster at $d = 3$.}
\emph{Error.} Leech rank $24 > h^{1,1}(K3) = 20$.
\emph{Ghost.} Fake Monster at $d = 5$ on $K3 \times K3 \times E$ via
$E_5$-Poisson shift-law row (Theorem~\ref{wn:thm:spine-dimension-census}).

\item \textbf{$\chi_{\cV_{24}}$ directly matches $\Delta_5^{-2}$.}
\emph{Error.} Virasoro $(2, 45)$-minimal apparatus inapplicable.
\emph{Ghost.} Heisenberg--Mukai identity
$\eta^{-48} = (2\pi i z)^2\,\Delta_5^{-2}|_{H_1, z = 0, \tau = \sigma}$
at the diagonal of the Humbert divisor
(Theorem~\ref{wn:thm:spine-heis-mukai}).

\item \textbf{Gaiotto curve $\Sigma_{2,0}$ with $c_{4d} = 107/6$.}
\emph{Error.} Genus-$2$ closed curve gives $c_{4d} \in \{7/6, 13/6\}$;
the $+90/12$ "Beem--Rastelli Coulomb regulator" is a phantom citation
--- Beem--Rastelli $2015$ \S 4.3 has no such regulator shift.
\emph{Ghost.} Class-$\mathcal{S}$ datum $(A_1, \Sigma_{0,24})$
($24$-punctured sphere) via Chacaltana--Distler $2010$ Table~$3$
row~$1$: $c_{4d} = (5n - 13)/6 = 107/6$ at $n = 24$, conjectural
pending primary-source trace of the nilpotent-Higgsing combinatorics.

\item \textbf{$\Phi$ natively produces an $E_n$-chiral algebra on a
curve.}
\emph{Error.} Backwards framing.
\emph{Ghost.} Two-stage factorisation
(Theorem~\ref{wn:thm:spine-two-stage}): $\Phi^{\mathrm{FA}}_d$
produces the canonical $E_d^{\mathrm{hol}}$-hFA on $X$;
$\mathrm{Sp}^{\mathrm{ch}}_{\Sigma_{d-1}, C}$ specialises to a
curve.

\item \textbf{Shifted-symplectic table terminates at $d = 4$.}
\emph{Error.} Wrong.
\emph{Ghost.} PTVV admits arbitrary integer shifts; at $d = 5$,
$E_5$-Poisson with shift $+1$
(Theorem~\ref{wn:thm:spine-hCS-classical}).

\item \textbf{Trace of Hecke $T_p$ on $S_5(\Spr_4(\ZZ),
\nu_{\Delta_5})$ equals $1$.}
\emph{Error.} Multiplicity-one ($\dim = 1$) conflated with Hecke
eigenvalue ($\lambda_p$, transcendental at generic $p$).
\emph{Ghost.} $\lambda_p(\Delta_5) = \chi_{\mathrm{spin}}(p)
\sqrt{\lambda_p(\Delta_{10})}$ metaplectic square-root form.

\item \textbf{Twisted-twined $H^3(C_{\Mtf}(g_N); U(1)) = 0$ uniformly.}
\emph{Error.} Non-uniform across the nine admissible cells.
\emph{Ghost.} Seven cells ($3B, 4A, 4B, 4C, 6A, 6B$, and one more)
vanish cleanly; $2A$ carries $\ZZ/2$, $2B$ carries $\ZZ/4$ residual
discrete torsion (Milgram $1995$; Benson $1998$). $\kBKM(\Phi_2) =
c_2(0)/2 = 4$ is cocycle-free; torsion dresses the paramodular
multiplier.

\item \textbf{$\CoHA(\CC^3) = Y^+$ strict.}
\emph{Error.} Requires base change to the localisation
$\mathbb{F} = \CC(\epsilon_1, \epsilon_2)$.
\emph{Ghost.} Three-presentation theorem at localisation
(Theorem~\ref{wn:thm:spine-coha-miki}); integer-unlocalised version
has cokernel killed by $\epsilon_1\epsilon_2\epsilon_3$.

\item \textbf{$\Winf$-triality is source of $Y^+$-triality.}
\emph{Error.} Direction backwards.
\emph{Ghost.} Single source is shuffle-kernel $S_3$-symmetry;
$\Winf[\lambda]$-triality is the $\mathrm{ev}_\lambda$-image shadow
of $Y^+$-triality.

\item \textbf{CY$_3$ "gains" $S_3$-triality.}
\emph{Error.} Direction backwards.
\emph{Ghost.} The $S_3$-triality is ambient on the three-parameter
$(\epsilon_1, \epsilon_2, \epsilon_3)$ space; the CY slice is a
codimension-$1$ hyperplane that \emph{restricts} the triality
faithfully, not creates it.

\item \textbf{$\CoHA(\CC^3)$-Yangian acts on $H^*_T(\mathrm{Hilb}^n
(\CC^3))$.}
\emph{Error.} $\mathrm{Hilb}^n(\CC^3)$ is non-smooth for $n \geq 4$
(Iarrobino $1972$) and non-irreducible for $n \geq 8$
(Brian\c con $1977$); no Nakajima Heisenberg action (no holomorphic
$2$-form on $\CC^3$).
\emph{Ghost.} The correct module is
$\bigoplus_n H^*_T(\mathrm{Hilb}^n(\CC^2))$ with $\epsilon_3$ entering
through the shuffle kernel, equivalently
$\bigoplus_n H^*_T(\cM_n(\CC^3), \phi_W)$ (Donaldson--Thomas cohomology
of the Jordan triple quiver moduli).

\item \textbf{Prochazka--Rapcak embedding
$Y(\widehat{\mathfrak{sl}}_n) \hookrightarrow Y(\widehat{\fgl}_1)$.}
\emph{Error.} Direction reversed.
\emph{Ghost.} Prochazka--Rapcak $2018$ Thm.~1.1 gives a
\emph{truncation} $\cW_\infty[\lambda] \twoheadrightarrow \cW_n$
with level rescaling; the Yangian-generator embedding is a separate
theorem.

\item \textbf{Commutativity of $E_3^{\mathrm{hol}}$ from
$\pi_1(S^5) = 0$.}
\emph{Error.} $\pi_1(\mathrm{Conf}_2(\CC^3)) = 0$ gives only
pair-exchange symmetry; full $E_3$-structure requires higher homotopy
data.
\emph{Ghost.} Genuine $E_3^{\mathrm{hol}}$-braiding sources from
$\bar\partial$-grading of Dolbeault chains plus $\Omega$-background
spectral parameters, not from $\pi_1(S^5)$.

\item \textbf{$\overline{\mathrm{Conf}}_n(\CC^3)$ as the operadic
base.}
\emph{Error.} Bochner--Martinelli propagator has conditional
convergence on $\overline{\mathrm{Conf}}_n$; the operadic composition
does not converge absolutely.
\emph{Ghost.} The correct base is the Axelrod--Singer /
Fulton--MacPherson compactification $\FMcpt(n)$ with blowup corner
charts; convergence is absolute there.

\item \textbf{GW21 strict $\leftrightarrow$ FG12 homotopy $E_3$-Koszul
is a "follows from".}
\emph{Error.} Composition is not automatic.
\emph{Ghost.} Three-step chain: Gwilliam--Williams strict, composed
with Fresse $2017$ Vol.~I Thm.~12.3.A (bar--cobar Quillen equivalence),
composed with Positselski $2011$ coderived/contraderived transfer.

\item \textbf{"Propagator kills harmonic subspace" (minimal model).}
\emph{Error.} $P_{\mathrm{BM}}$ is a Green operator, not an
annihilation operator.
\emph{Ghost.} Dolbeault degree count via Hodge homotopy: harmonic
inputs sit in $\Omega^{0,0}$; Hodge propagator maps
$\Omega^{0,q} \to \Omega^{0,q-1}$; target $\Omega^{0,-1}$ vanishes
(Kapranov $1999$ \S 3.1).

\item \textbf{$\mathrm{At}(TX) = 0$ is the formality obstruction on
compact CY$_3$.}
\emph{Error.} $\mathrm{At}(TX) = 0$ is necessary, not sufficient.
\emph{Ghost.} Full Kapranov tower $\{\kappa_n\}_{n \geq 3}$; leading
$\kappa_3 = \mathrm{At}\cup\mathrm{At}$; higher coefficients are
Duflo-class contributions (Markarian $2009$).

\item \textbf{Kuranishi cubic lives in $H^3(K3, \Omega^3_{K3})$.}
\emph{Error.} Wrong receptacle.
\emph{Ghost.} Correct receptacle is $H^3(K3, \Lambda^3 T_{K3})$, which
vanishes by rank ($T_{K3}$ rank $2$, $\Lambda^3 = 0$).

\item \textbf{CY slice reduces $\mathrm{Sym}^2(\fg^\vee)^{\fg}$ from
rank-$2$ to rank-$1$.}
\emph{Error.} Wrong side.
\emph{Ghost.} $\mathrm{Sym}^2(\fg^\vee)^{\fg} = \CC\cdot B$ is
rank-$1$ for simple $\fg$ by Whitehead, independent of
$\epsilon_i$'s; the $\epsilon$-reduction happens on the Yangian side.

\item \textbf{$\mathrm{Sp}_4(\ZZ)/\{\pm I\} \simeq \SOr(\Lambda^{3,2})$
via $\wedge^2$-Pfaffian.}
\emph{Error.} The exterior-square lands in $\mathrm{SO}$ of a
rank-$6$ lattice of signature $(3,3)$; descent to $(3,2)$ requires
Pl\"ucker $\mathrm{LG}(2,4) \hookrightarrow \bP(\wedge^2 V_4)$.
\emph{Ghost.} Finite-index embedding
$\Spr_4(\ZZ)/\{\pm I\} \hookrightarrow \SOr(\Lambda^{3,2})$ with
non-trivial cokernel (Maass multiplier class).

\item \textbf{Eight-form Borcherds-weight sequence
$\{5, 2, 1, 1, 1/2, 1, 1/4, 0\}$ extends the CHL Borcherds ladder
$\{5, 2, 1, 1, 1\}$.}
\emph{Error.} Conflates two distinct constructions on different
Jacobi inputs.
\emph{Ghost.} $\{5, 2, 1, 1, 1\}$ is the CHL Borcherds-weight
refinement of $\phi_{0,1}$; $\{5, 4, 3, 2, 1\}$ is the Gritsenko
additive-lift of weight-$k(N)$ index-$1$ Jacobi forms; the two
agree only at $N = 1$ by accidental equality
$\phi_{0,1} = \mathrm{Grit}(\phi_{0,1})$. The half-integer rows
at $N \in \{5, 7, 8\}$ are Weil-rep central-extension indices via
Shimura--Waldspurger.

\item \textbf{Borcherds Monster lives on $(T^{24}_{\LLeech}\times E)/
\ZZ_2$ as a CY-$3$.}
\emph{Error.} $\dim_\CC((T^{24}_{\LLeech}\times E)/\ZZ_2) = 25$, not
$3$.
\emph{Ghost.} Borcherds Monster $\fm$ on $\mathrm{II}_{1,1}$, virtual
CY-$3$ $\Psi$-row. FLM orbifolds the lattice VOA $V_{\LLeech}$, not
the lattice variety. The $V^\natural$ construction is chiral-boundary
shadow of a virtual CY-$3$ datum; it is not a genuine Stage-$2$
specialisation of $\Phi^{\mathrm{FA}}_3$.

\item \textbf{$\fg_{\Delta_5}$ is the $\sigma_q$-fixed subalgebra of
$\fg^{\mathrm{BKM}}_{\Lambda^{3,3}}$.}
\emph{Error.} Borcherds $1998$ Thm.~13.3 requires $b^- = 2$; at
signature $(3, 3)$ the Grassmannian is real and carries no
holomorphic $\mathrm{O}^+$-form. The three BKMs $\fm$, $\fg_{\mathrm{FM}}$,
$\fg_{\Delta_5}$ have Cartans of ranks $26, 26, 3$, pairwise
non-embeddable.
\emph{Ghost.} Hyperbolic restriction correspondence: $\Phi_{12}$ on
$\mathrm{II}_{26,2}$ restricts along $\mathrm{II}_{2,2}
\hookrightarrow \mathrm{II}_{25,1}$ to $\Phi_{10} = \Delta_5^2$
(Borcherds $1998$ \S 14 K\"unneth).

\item \textbf{Hauptmodul from Stage-$1$ uniqueness up to contractible
choice.}
\emph{Error.} Conflates two independent arithmetic inputs.
\emph{Ghost.} Hauptmodul of McKay--Thompson series is Borcherds $1992$
Hecke-replication on the Monster root lattice, independent of the
$\Phi^{\mathrm{FA}}_3$-rigidity statement.

\item \textbf{Six routes to $G(K3 \times E)$ as six
$(\Sigma_2, C)$-specialisations.}
\emph{Error.} Only three routes admit cycle-class indexing.
\emph{Ghost.} Three $(\Sigma_2, C)$-specialisations plus three routes
consuming non-cycle-class data (a Jacobi form, a lattice, a 6d
SCFT); correct indexing is the triple (orbit of $(\Sigma_2, C)$,
construction machine, generator-rank $\rho^{R_i} \in \{3, 12, 24\}$).

\item \textbf{$K^{\kappa_{\mathrm{ch}}}_\mathsf{B} =
\kch(\cH_{\Mukr}) + \kBKM(\Phi_1) = 3 + 5 = 8$.}
\emph{Error.} $\kch(\cH_{\Mukr}) = 4$, not $3$; $\kch^!$ and
$\kBKM$ are different invariants.
\emph{Ghost.} $K^{\kch}_\mathsf{B} = \varrho \cdot K = (1/6)\cdot 48
= 8 = 2c_+(\Mukr(K3))$; three faces of $8$ unified via Bruinier
reciprocity and Drinfeld scaling
(Theorem~\ref{wn:thm:spine-five-archetype}).

\item \textbf{$c_N = 24 k_N$ reduced central charge.}
\emph{Error.} At $N = 1$: $c_1 = 24 \cdot 5 = 120$, contradicting
MSW $c_L = 6k + 24$.
\emph{Ghost.} $c_L^{\mathrm{reduced}} = 24$ universal in $N$; $k_N$
is the \emph{Borcherds weight} of $\Phi_{k_N}$, not a central charge.

\item \textbf{$h^{2,2}(K3 \times K3) = 402$.}
\emph{Error.} Missed two corner contributions $h^{0,0}\cdot h^{2,2}
+ h^{2,2}\cdot h^{0,0} = 2$.
\emph{Ghost.} $h^{2,2}(K3\times K3) = 1 + 1 + 400 + 1 + 1 = 404$.

\item \textbf{Naive affine folding $\widehat{A_{2n-1}}^\sigma
\to \widehat{B_n}$.}
\emph{Error.} Untwisted vs twisted affine distinction.
\emph{Ghost.} Correct folding is twisted affine
$\widehat{A_{2n-1}}^\sigma = A_{2n-1}^{(2)}$ with shifted grading
(Kac Aff~2, Aff~3); the twisted affine algebra, not the untwisted
$\widehat{B_n}$.

\item \textbf{Chan--Paton minuscule works for all ADE.}
\emph{Error.} Fails for $E_8$ (no minuscule representation).
\emph{Ghost.} Chan--Paton minuscule clean for $A_n, D_n, E_6, E_7$;
Kostant--Slodowy slice for $E_8$ with output
$\cW_{\mathrm{fin}}(\fe_8) \subsetneq Y(\widehat{\fe}_8)$. The two
are different constructions giving different outputs.

\item \textbf{Costello--Dimofte--Paquette $2020$ rank-one
all-orders base.}
\emph{Error.} Phantom citation.
\emph{Ghost.} Actual rank-one base is Costello $2013$ \S 10 +
Costello--Gaiotto $2018$ \S 6.

\item \textbf{RSYZ $2020$ cohomological toroidal
$\widehat{\widehat{\fgl}_2}$.}
\emph{Error.} Miscited.
\emph{Ghost.} Correct attribution: Negu\c t $2015$ (K-theoretic
shuffle); Feigin--Jimbo--Miwa--Mukhin $2016$ (toroidal shuffle);
Kapranov--Vasserot $2018$ (geometric K-theoretic action).
\end{enumerate}
\end{theorem}

\subsection{Residual frontier}
\label{wn:subsec:spine-frontier}

\begin{itemize}
\item 3-dualizability on general compact CY$_3$
(Conjecture~\ref{wn:conj:spine-compact-recovery}): requires an
$E_3$-coherent-duality extension of Gwilliam--Williams $2021$
Prop.~5.3.2.
\item $(\infty,1)$-functoriality of $\Phi^{\mathrm{FA}}_d$ on morphisms
at $d \geq 3$ (Pattern 273): open even rationally when the CY
category is non-formal (local~$\bP^2$).
\item Integral $E_d$-formality at $d \geq 3$: the integral identity of
topological $E_d$ and algebraic $E_d$ via the Gerstenhaber bracket of
degree $1-d$ is open at $d \geq 3$ (conceded in
\texttt{chapters/theory/quantum\_chiral\_algebras.tex} line~1033).
\item BCFG extension of the all-orders compound theorem: bottleneck is
the $\sigma$-equivariant renormalisation scheme for the Costello 6D
$\hCS$ construction under the Dynkin involution.
\item Elliptic-surface specialisation $(\Sigma_2, C) = (\mathcal{E},
\bP^1)$ with Mordell--Weil-indexed simple roots: conjectural for K3 of
Picard rank $20$ (Shioda--Inose); scope hypothesis
$\rho(K3) = 20$ is necessary.
\item Non-abelian Fake Monster at $d = 5$:
$\chi_{A^{\mathrm{FM}}_E}(q, Z) = 1/\Phi_{12}(Z)$ after Niemeier
projection, conjectural; doubly-reduced Donaldson--Thomas integrand
analogue of Oberdieck--Pixton on $K3\times K3\times E$.
\item Non-CHL $N = 7$: order-$4$ central extension of $\Mpr_4$ by
$\mu_4$ (not a Chern--Simons gerbe), open.
\item Class-$\mathcal{S}$ $(A_1, \Sigma_{0,24})$ with $c_{4d} = 107/6$:
conjectural, falsifiable at $q^{10}$ via direct Schur-index
computation (Gadde--Rastelli--Razamat--Yan $2013$ TQFT prescription).
\item Ran-level $S_3$-triality of Miki as a factorisation-algebra
automorphism: open.
\item Rank-$\geq 3$ lattice-polarised $\fg_L$ family: extension of the
envelope $\fg_{\Lambda^{3,2}}$ to higher-rank Picard lattices.
\item Bracket-level $\fg_{\mathrm{BPS}}(K3 \times E) \simeq
\fg_{\Delta_5}$: reduces to a single Hecke--Borcherds identity on the
Gritsenko $1999$ paramodular family, verifiable via Harvey--Moore
one-loop amplitudes.
\end{itemize}

\subsection{The picture in one sentence}
\label{wn:subsec:spine-one-line}

\emph{A single $E_d^{\mathrm{hol}}$-holomorphic factorisation algebra
$\cF_X$ on each Calabi--Yau $d$-fold $X$ is the canonical Stage-$1$
output of $\Phi$, realised at $d = 3$ as the
Costello--Francis--Gwilliam $2026$ $E_3^{\mathrm{hol}}$-algebra of
observables of six-dimensional holomorphic Chern--Simons on $X$;
transverse-bordism classes of $(\Sigma_{d-1}, C)$-pairs specialise
$\cF_X$ to $E_1$-chiral shadows --- Igusa $\Delta_5$ GBKM, Borcherds
Monster, Fake Monster (at $d = 5$), Niemeier cousins, CHL descendants
--- and the universal Borcherds-weight identity
$\kBKM(\Phi_N) = c_N(0)/2$ holds at two scopes, the CHL
Borcherds-weight refinement of $\phi_{0,1}$ giving $\{5,2,1,1,1\}$ and
the Gritsenko additive-lift of weight-$k(N)$ Jacobi inputs giving
$\{5,4,3,2,1\}$; the six routes to the quantum vertex chiral group
$G(K3\times E)$ are three $(\Sigma_2, C)$-specialisations of
$\cF_{K3\times E}$ plus three routes consuming non-cycle-class data,
not six applications of $\Phi$, and they are indexed by the triple
$(\text{$(\Sigma_2, C)$-orbit},\; \text{construction machine},\;
\text{generator-rank } \rho^{R_i})$ with $\rho^{R_i} \in \{3, 12, 24\}$.}

% End of post-adversarial platonic synthesis.
